\section{Gaussian Space --- Noise --- Exponential}
%%% generated by the blueprint pipeline from TCSlib.BooleanAnalysis.GaussianSpace.Noise.Exponential
\subsection*{Gaussian Space --- Noise --- Exponential: Introduction}

Declarations of the Lean module \texttt{TCSlib.BooleanAnalysis.GaussianSpace.Noise.Exponential}.

\subsection*{Gaussian Space --- Noise --- Exponential: Results}

\begin{lemma}[Exponential moments of Gaussian coordinate sums]
\lean{BooleanAnalysis.GaussianSpace.integrable_exp_sum_abs_gaussian}
\label{BooleanAnalysis.GaussianSpace.integrable_exp_sum_abs_gaussian}
\leanok
\uses{BooleanAnalysis.ThresholdFunctions.standardGaussianMeasure}
\difficulty{4}
For every natural number $n$ and every real number $c$, the function
$x\mapsto\exp\!\left(c\sum_{i=1}^n |x_i|\right)$ is integrable with respect to the
product of $n$ standard Gaussian measures on $\mathbb R^n$.
\end{lemma}

\begin{lemma}[Exponential of Gaussian coordinate magnitudes is integrable]
\lean{BooleanAnalysis.GaussianSpace.memLp_exp_sum_abs_gaussian}
\label{BooleanAnalysis.GaussianSpace.memLp_exp_sum_abs_gaussian}
\leanok
\uses{BooleanAnalysis.GaussianSpace.integrable_exp_sum_abs_gaussian}
\difficulty{4}
Let $n$ be a natural number, let $c\in\mathbb R$, and let $p\in[0,\infty]$ be finite. On
$\mathbb R^n$ with product standard Gaussian measure, the function $x\mapsto
\exp\!\left(c\sum_{i=1}^n |x_i|\right)$ belongs to $L^p$.
\end{lemma}

\begin{lemma}[Dominated convergence in $L^p$]
\lean{BooleanAnalysis.GaussianSpace.tendsto_eLpNorm_zero_of_dominated}
\label{BooleanAnalysis.GaussianSpace.tendsto_eLpNorm_zero_of_dominated}
\leanok
\difficulty{6}
Let $(\Omega,\mu)$ be a measure space, let $E$ be a normed additive group, and let $p$
be a nonzero finite extended nonnegative real number. Suppose $(F_k)_{k\in\mathbb N}$
consists of functions $\Omega\to E$ that are strongly measurable almost everywhere, and
that a real-valued function $b$ belongs to $L^p(\mu)$. If $\|F_k(x)\|\le |b(x)|$ for
almost every $x$ for each $k$, and $F_k(x)\to0$ for almost every $x$, then
$\|F_k\|_{L^p(\mu)}\to0$.
\end{lemma}

\begin{lemma}[Taylor approximation of the Gaussian exponential]
\lean{BooleanAnalysis.GaussianSpace.tendsto_gaussian_exp_taylor}
\label{BooleanAnalysis.GaussianSpace.tendsto_gaussian_exp_taylor}
\leanok
\uses{BooleanAnalysis.GaussianSpace.memLp_exp_sum_abs_gaussian, BooleanAnalysis.GaussianSpace.tendsto_eLpNorm_zero_of_dominated}
\difficulty{5}
Let $n\ge 0$, let $L:\mathbb R^n\to\mathbb C$ be a continuous real-linear map, and let
$p\in[0,\infty]$ satisfy $p<\infty$. Under standard Gaussian measure $\gamma_n$ on
$\mathbb R^n$, the $L^p$ quasi-norm of the error between $e^{L(x)}$ and its
degree-$(k-1)$ Taylor polynomial tends to zero as $k\to\infty$; explicitly,
$\|e^{L(\cdot)}-\sum_{m=0}^{k-1}L(\cdot)^m/m!\|_{L^p(\gamma_n)}\to0$.
\end{lemma}
